\section{Methodology}
\label{sec:methodology}

\subsection{Primary Climate-Scene Dataset and Provenance Reconstruction}

The primary dataset contains 6,334 aerial and remote-sensing images
assigned to three semantic classes: 1,931 \emph{Desert}, 2,338
\emph{Flooded}, and 2,065 \emph{Neither}. The dataset is a heterogeneous
collection assembled from multiple underlying image sources rather than
a single controlled acquisition campaign. Consequently, images belonging
to the same semantic class can differ substantially in acquisition
platform, geographic context, spatial resolution, visual style, and
surrounding scene composition.

To support source-aware evaluation, image-level provenance was
reconstructed for the complete dataset. Eight source groups were
identified, and every image was assigned to one source group with no
unresolved records. This reconstruction is essential because conventional
random splitting can place visually related images originating from the
same source family into both training and test partitions. In contrast,
source-disjoint evaluation can hold entire source groups out of model
development.

The source labels are interpreted as provenance groups rather than
perfectly controlled imaging modalities. Source identity, semantic class,
sensor characteristics, acquisition conditions, and platform are partially
entangled within the assembled dataset. Therefore, the source-disjoint
experiments are treated as composite source-distribution shifts and not
as pure causal UAV-versus-satellite modality experiments.

\subsection{Input Processing and Baseline Augmentation}

All models receive RGB inputs resized to \(224\times224\) pixels. Images
are converted to floating-point values in the interval \([0,1]\). The
class-index mapping is fixed throughout all experiments as

\[
[\text{Desert},\text{Flooded},\text{Neither}]
\rightarrow
[0,1,2].
\]

Augmentation is applied only to the training partition. Validation,
internal test, and external test images are not randomly augmented. The
baseline augmentation pipeline consists of horizontal-and-vertical random
flipping, random rotation with factor \(0.08\), random zoom with factor
\(0.10\), and random contrast with factor \(0.10\). The same baseline
augmentation family is used for CustomCNN\_V1 and the MobileNetV2 ERM
baseline so that architectural comparisons are not confounded by different
standard augmentation policies.

Random seeds are explicitly controlled. The principal baseline
configuration uses seed 42, and TensorFlow deterministic operations are
enabled where supported. Training batches contain 16 images.

\subsection{CustomCNN\_V1 Baseline}

A compact convolutional network, denoted CustomCNN\_V1, is trained from
scratch to provide a non-pretrained reference model. Its purpose is not to
maximize architectural complexity but to establish how far a comparatively
small task-specific CNN can generalize when no external representation
pretraining is used.

The network contains three convolutional feature-extraction blocks. The
first block applies a \(3\times3\) convolution with 32 filters and
``same'' padding, followed by Batch Normalization, ReLU activation, and
max pooling. The second block repeats the same structure with 64 filters,
and the third block uses 128 filters:

\[
32 \rightarrow 64 \rightarrow 128.
\]

After the final pooling operation, Global Average Pooling converts the
spatial feature tensor into a fixed-length representation. This
representation is passed to a fully connected layer with 128 ReLU units,
followed by dropout with probability \(0.40\). The classifier terminates
in a three-unit softmax layer corresponding to Desert, Flooded, and
Neither.

The architecture contains 111,043 total parameters, of which 110,595 are
trainable and 448 are non-trainable Batch Normalization parameters. The
model is optimized with Adam using an initial learning rate of \(10^{-3}\)
and Sparse Categorical Cross-Entropy loss. Model selection is based
exclusively on validation loss. The lowest-validation-loss checkpoint is
retained, Early Stopping uses patience 5 with restoration of the best
weights, and the learning rate is reduced by a factor of \(0.5\) after
two validation-loss plateau epochs, with a minimum learning rate of
\(10^{-6}\). The internal test partition is not used for checkpoint
selection or training decisions.

\subsection{MobileNetV2 Transfer-Learning Baseline}

The principal transfer-learning model uses MobileNetV2 initialized from
ImageNet weights \cite{sandler2018mobilenetv2}. The convolutional
backbone is instantiated without its original ImageNet classification
head. The same baseline augmentation used by CustomCNN\_V1 is first
applied to the normalized \([0,1]\) inputs. Because MobileNetV2 expects
the ImageNet preprocessing range \([-1,1]\), the augmented images are
then transformed according to

\[
x_{\mathrm{MN}} = 2x - 1.
\]

The pretrained backbone produces the learned image representation. The
task-specific classifier consists of Global Average Pooling, dropout with
rate \(0.30\), and a three-unit softmax output layer.

Training is performed in two stages. During Stage 1, the pretrained
backbone is frozen and only the new classification head is optimized using
Adam with learning rate \(10^{-3}\). During Stage 2, the final 30 backbone
layers are made eligible for fine-tuning while Batch Normalization layers
remain frozen; optimization then continues with learning rate \(10^{-5}\).
Validation loss determines checkpoint selection throughout the
transfer-learning procedure, and test samples remain isolated from model
selection.

This two-stage configuration allows the newly initialized classifier to
adapt first while preserving the pretrained representation, followed by
conservative adaptation of higher-level MobileNetV2 features to the
climate-scene task.

\subsection{Domain-Generalization Interventions}

The most difficult source-disjoint setting is subsequently used to
investigate whether training interventions can improve robustness beyond
the MobileNetV2 ERM baseline. Three intervention families are considered.

\subsubsection{Conservative Augmentation (DG-AUG-V2)}

DG-AUG-V2 retains the complete baseline augmentation pipeline and
introduces two additional transformations intended to increase
acquisition-style variability without severely changing scene semantics.
Random Translation uses height and width factors of \(0.05\) with
reflected boundary filling. Random Brightness uses factor \(0.08\) with
input range \([0,1]\).

Thus, the DG-AUG-V2 pipeline extends

\[
\{\text{flip},\text{rotation},\text{zoom},\text{contrast}\}
\]

with

\[
\{\text{translation}_{0.05},\text{brightness}_{0.08}\}.
\]

An earlier, more aggressive diagnostic configuration included stronger
perturbations and saturation changes and produced severe image
degradation. It is therefore not treated as a primary DG result.
DG-AUG-V2 is the frozen conservative augmentation configuration used in
the reproducibility analysis.

\subsubsection{MixStyle Diagnostics}

MixStyle perturbs feature statistics by mixing instance-level style
information between samples, with the intention of reducing dependence on
domain-specific appearance statistics \cite{zhou2021mixstyle}. Two
insertion variants were evaluated diagnostically at different positions
within the MobileNetV2 feature hierarchy. These experiments are retained
as diagnostic evidence but are not included in the three-seed comparison
because they were not replicated over the complete seed set.
Consequently, no claim of seed-stable MixStyle improvement is made.

\subsubsection{Representation Self-Challenging}

Representation Self-Challenging (RSC) is designed to reduce dependence on
the most immediately predictive features by suppressing dominant feature
responses during training \cite{huang2020rsc}. The implementation uses a
custom training step built around the MobileNetV2 representation. The
frozen configuration uses a feature-drop fraction of \(1/3\) and a
batch-challenge fraction of \(1/3\). The remaining model structure,
baseline augmentation, class mapping, and transfer-learning framework
follow the MobileNetV2 baseline as closely as possible so that performance
differences can be interpreted primarily in terms of the RSC
intervention.

Because DG performance may vary substantially with stochastic
optimization, the primary robustness comparison is not based on a single
run. ERM, DG-AUG-V2, and RSC are each repeated using seeds 42, 123, and
2026 under the same frozen source-disjoint protocol.

\subsection{External FloodNet Evaluation}

Internal source-disjoint testing remains limited to the source groups
available inside the primary dataset. External validation therefore uses
the independently collected FloodNet UAV dataset
\cite{rahnemoonfar2021floodnet}.

The official labeled FloodNet classification subset contains 398 images:
51 labeled \emph{Flooded} and 347 labeled \emph{Non-Flooded}. For
compatibility with the primary taxonomy, Flooded maps directly to the
internal Flooded class, while Non-Flooded maps to Neither. FloodNet
contains no supported Desert target class.

Importantly, the trained classifier is not converted into a two-class
model during external evaluation. Its original output vector

\[
[p_D,p_F,p_N]
\]

is preserved. If the model predicts Desert for a FloodNet sample, that
prediction remains an error rather than being discarded and the remaining
probabilities renormalized. External Macro-F1 is therefore reported over
the two supported target classes, Flooded and Neither, while unsupported
Desert predictions continue to affect the confusion structure and
accuracy.

Before evaluation, image-level leakage is screened against the complete
primary dataset. Exact duplicate screening uses SHA-256 hashes, while
perceptual hashes are used to search for visually near-duplicate images
that may survive resizing or re-encoding. This audit is performed before
interpreting external performance as independent evidence.

\subsection{Grad-CAM Explainability Protocol}

Grad-CAM is used to visualize spatial regions contributing to selected
class predictions \cite{selvaraju2017gradcam}. Explainability analysis is
applied only after model training and evaluation are frozen.

External FloodNet images are divided according to agreement across all
nine primary externally evaluated models: three methods---ERM,
DG-AUG-V2, and RSC---each trained with three seeds. From these predictions,
three behavior groups are defined:

\begin{enumerate}
    \item consensus-correct Flooded,
    \item consensus-correct Neither, and
    \item consensus Neither-to-Flooded failure.
\end{enumerate}

Within each behavior group, low-, median-, and high-confidence
representative examples are selected algorithmically before heatmap
interpretation. Grad-CAM is then generated for the seed-42 ERM,
DG-AUG-V2, and RSC models, yielding \(3\times3\times3=27\) heatmaps.

Grad-CAM is treated as qualitative explanatory evidence rather than a
causal description of the model decision process. Heatmap interpretation
is therefore complemented by a separate semantic analysis using FloodNet
segmentation masks.

\subsection{Segmentation-Based Contextual Shortcut Audit}

To investigate whether recurring external false-flood predictions are
associated with particular semantic scene elements, official FloodNet
segmentation masks are analyzed for true Neither samples.

Two consensus cohorts are defined using all nine externally evaluated
models. A \emph{consensus failure} is a true Neither image classified as
Flooded by all nine models, whereas a \emph{consensus correct} case is
classified as Neither by all nine models.

For each available segmentation class, two complementary quantities are
computed. First, class presence records whether at least one pixel of that
semantic category is present in an image. Second, pixel coverage records
the fraction of the image occupied by that category. Presence-rate
differences between failure and correct cohorts are evaluated using
Fisher's exact test. Pixel-coverage distributions are compared using the
Mann--Whitney \(U\) test because the class fractions are not assumed to
follow a Gaussian distribution.

Because multiple semantic classes are examined simultaneously,
\(p\)-values are adjusted using the Holm procedure across the ten
segmentation categories. This analysis is explicitly post hoc and
exploratory. A statistically significant association between a scene
element and model failure indicates that the element is
disproportionately represented in the failure cohort; it does not by
itself establish that the element causally produced the model prediction.

\subsection{Evaluation Metrics}

Macro-F1 is the primary classification metric because it weights each
semantic class equally despite differences in class frequency. For class
\(c\),

\[
P_c =
\frac{TP_c}{TP_c + FP_c},
\qquad
R_c =
\frac{TP_c}{TP_c + FN_c},
\]

and

\[
F1_c =
\frac{2P_cR_c}{P_c + R_c}.
\]

For \(C\) supported classes, Macro-F1 is

\[
F1_{\mathrm{macro}}
=
\frac{1}{C}
\sum_{c=1}^{C} F1_c.
\]

For the internal three-class experiments, \(C=3\). For the
supported-class FloodNet metric, \(C=2\) for Flooded and Neither, while
Desert predictions remain errors in the original three-class prediction
process.

Secondary evaluation measures include overall accuracy, class-wise
precision and recall, confusion matrices, one-versus-rest area under the
receiver-operating-characteristic curve (AUROC), and area under the
precision--recall curve (AUPRC). Multi-seed results are summarized using
the arithmetic mean and sample standard deviation across seeds.
